\chapter{1931:Carl Bosch and Friedrich Bergius}

\label{chap:chemistry1931}

The Nobel Prize in Chemistry for the year 1931 was awarded jointly to Carl Bosch and Friedrich Bergius.

\textbf{Motivation:}

in recognition of their contributions to the invention and development of chemical high pressure methods

\section{Carl Bosch}

\subsection*{Early Life and Education}

Carl Bosch was born on 1874-08-27 in Cologne, Germany.
He was a German chemist and industrialist who shared the 1931 Nobel Prize in Chemistry with Friedrich Bergius for their contributions to the invention and development of chemical high pressure methods.

\subsection*{Career and Major Research}

Bosch studied chemistry at the Technische Hochschule in Charlottenburg and at the University of Leipzig, where he received his doctorate in 1898.
He joined the Badische Anilin- und Soda-Fabrik (BASF) in 1899 and devoted himself to industrial chemistry.
From 1909 he led the task of turning Fritz Haber's laboratory synthesis of ammonia into an industrial process, developing the high-pressure reactors, the catalysts, and the large-scale production of hydrogen that the process required.
He also worked on the hydrogenation of coal to liquid fuels and on the synthesis of methanol under high pressure.
After the formation of IG Farben in 1925 he served as chairman of its management board, and in 1933 he succeeded Fritz Haber as head of the Kaiser Wilhelm Institute for Physical Chemistry in Berlin-Dahlem.
He protested publicly against the Nazi dismissal of Jewish scientists, warning that it would destroy German science.

\subsection*{Nobel Prize-winning Work}

The work for which Carl Bosch was awarded the Nobel Prize in Chemistry in 1931 was described by the Nobel Committee as: "in recognition of their contributions to the invention and development of chemical high pressure methods".

Bosch's contribution was the development of the chemical high pressure methods by which industrial synthesis under extreme pressures and temperatures became practical.
His work made possible the large-scale production of ammonia from atmospheric nitrogen, and thereby the production of nitrogen fertilizers.

\subsection*{Significance and Impact}

The Bosch process for ammonia synthesis is among the most important chemical processes ever developed, because it provides the nitrogen compounds on which modern agriculture depends.
High pressure technology later found applications in the hydrogenation of coal and in the synthesis of fuels and chemicals.

\subsection*{Later Career and Honors}

Bosch directed the affairs of IG Farben until his death and remained one of the leading figures of the German chemical industry.
He was elected to the German Academy of Sciences Leopoldina and held honorary doctorates from several universities.
He died on 1940-04-26 in Heidelberg, Germany.

\subsection*{Legacy}

The legacy of Carl Bosch endures through the lasting impact of his scientific and technical contributions.
The industrial synthesis of ammonia that he made possible continues to supply the nitrogen fertilizers that sustain modern agriculture.

\subsection*{Modern Developments}

Bosch's industrial high-pressure chemistry, together with Haber, created the Haber–Bosch ammonia process that now produces roughly 200 million tons of ammonia per year, sustaining global food production. The process is today being transformed into green ammonia using electrolytic hydrogen from renewable energy, and high-pressure catalysis remains central to the modern chemical industry.

\subsection*{Selected Publications}

\begin{itemize}

\item Bosch, C. (1932). "The Development of the Chemical High Pressure Method during the Establishment of the New Ammonia Industry." Nobel Lecture, Stockholm.

\end{itemize}

\clearpage

\section{Friedrich Bergius}

\subsection*{Early Life and Education}

Friedrich Karl Rudolf Bergius was born on 1884-10-11 in Breslau, Germany (now Wrocław, Poland).
He was a German chemist who shared the 1931 Nobel Prize in Chemistry with Carl Bosch for their contributions to the invention and development of chemical high pressure methods.

\subsection*{Career and Major Research}

Bergius studied chemistry in Breslau and Leipzig, where he received his doctorate in 1907 under Arthur Hantzsch, working on reactions at high pressure.
After brief periods with Walther Nernst and Fritz Haber, he joined industrial research in 1911 and set up a laboratory at Rheinau near Mannheim for the study of coal.
There he developed the Bergius process, in which coal is hydrogenated with hydrogen under high pressure and temperature to give liquid fuels.
The process was developed industrially during and after the First World War, when Germany lacked access to petroleum.
After the Second World War he emigrated to the United States and later acted as a consultant in Argentina.

\subsection*{Nobel Prize-winning Work}

The work for which Friedrich Bergius was awarded the Nobel Prize in Chemistry in 1931 was described by the Nobel Committee as: "in recognition of their contributions to the invention and development of chemical high pressure methods".

Bergius's contribution was the invention and development of high pressure methods for the chemical transformation of coal.
He showed that bituminous coal can be hydrogenated under high pressure to give liquid hydrocarbons, providing a route from coal to liquid fuels.

\subsection*{Significance and Impact}

The Bergius process demonstrated the industrial power of high pressure chemistry, and the hydrogenation of coal and lignite came to supply a substantial part of Germany's fuel needs during the Second World War.

\subsection*{Later Career and Honors}

After the war Bergius worked as an industrial consultant in the United States and in Argentina, where he died on 1949-03-30 in Buenos Aires.
He held honorary doctorates from several universities.

\subsection*{Legacy}

The legacy of Friedrich Bergius endures through the lasting impact of his scientific contributions.
The high pressure hydrogenation of coal that he pioneered remains a reference point in the chemistry of fuels.

\subsection*{Modern Developments}

Bergius's high-pressure hydrogenation of coal and heavy oils laid the foundation for modern hydroprocessing in petroleum refining. His methods evolved into the catalytic hydrotreating and hydrocracking processes that produce clean transportation fuels. Bergius's vision is also relevant to current coal-to-liquid and biomass-to-liquid fuel technologies under development for energy security.

\subsection*{Selected Publications}

\begin{itemize}

\item Bergius, F. (1913). "Die Anwendung hoher Drucke bei chemischen Vorgängen und ein Nachbild des Entstehungsprozesses der Steinkohle." Halle: Wilhelm Knapp.

\item Bergius, F. (1932). "Chemical Reactions under High Pressure." Nobel Lecture, Stockholm.

\end{itemize}

\clearpage

\section*{Chapter Summary}

The 1931 prize recognized Carl Bosch and Friedrich Bergius for their contributions to the invention and development of chemical high-pressure methods. The industrial synthesis of ammonia and the hydrogenation of coal revolutionized chemical production.
